\chapter{Working efficiently}
\setwlchapter{workflow}
\section{Expand at the step that needs it}
For a four-fermion interaction written as $U D$ the compact notation
hides the number of terms obtained after a transformation of orbitals.
Replacing every creator by a sum of $M$ creators can produce $M^4$
terms before momentum selection. Apply Kronecker deltas while the
indices to be eliminated are still visible. A good intermediate
observable is often a commutator or a scalar matrix element rather
than the expanded Hamiltonian itself.
\casestudy{collect}{Separating expansion from collection of sums}
\casestudy{recognize}{Recovering the physical spin coupling}

\section{How large is the matrix?}
For $L$ spinful orbitals the unrestricted basis has $4^L$ states, so
a dense matrix contains $16^L$ entries. At $L=8$ a dense complex
double-precision array already requires 64 GiB just for its elements.
A sparse hopping Hamiltonian instead connects each Fock state to at
most a number of others proportional to the number of bonds. Symmetry
reduction and sparse storage address different parts of this cost.
For a spectral line one often needs only the initial sector and the
sectors connected by the probe, rather than all transitions.

\section{Tests before a long run}
Start with one bond. Check Hermiticity, a conserved quantity and
one matrix element whose sign is known. Next check a limiting value:
does the hopping spectrum reduce to on-site levels at $t=0$? Only
then increase the cluster. For degenerate eigenvalues compare spectral
projectors or subspace weights; two diagonalizers may choose different
orthonormal vectors inside the same eigenspace.

A short export example checks that the energy symbol and spin label
survive conversion to LaTeX:
\begin{wl}
snegfermionoperators[c];
snegrealconstants[\[Epsilon]];
tex = snegtex[\[Epsilon] number[c[], UP]];
StringContainsQ[tex, "\\epsilon"] && StringContainsQ[tex, "\\uparrow"]
\end{wl}
Save the mode list with any exported numerical matrix and record
whether additive constants and chemical-potential terms are included.
A sorted eigenvalue list alone cannot detect a reversed site order.
\begin{exercise}[An independent matrix test]\label{ex:trace}
Find the trace of $U\sum_i D_i$ over the unrestricted Fock space of
$L$ spinful orbitals. Would checking only that trace detect an
incorrect sign in the hopping? Give a second trace that would
distinguish missing bonds.
\end{exercise}
